\section{第6组：静电场}
\sourceproblem{6-T1}{23}{27}
\problemtext 半径 $R$ 的细带电圆环，线电荷密度随从正 $x$ 轴量起的极角变化为 $\lambda=\lambda_0\sin\theta$，求环心场强。
\workings $\mathrm dq=\lambda_0R\sin\theta\,\mathrm d\theta$，环心的场元为
\[
\mathrm d\bm E=-\frac{\lambda_0}{4\pi\varepsilon_0R}\sin\theta(\cos\theta\bm i+\sin\theta\bm j)\,\mathrm d\theta.
\]
在 $0$ 到 $2\pi$ 积分，交叉项为零，$\int\sin^2\theta\,\mathrm d\theta=\pi$，故 $\bm E=-\lambda_0\bm j/(4\varepsilon_0R)$。
\crosscheck 若 $\lambda_0>0$，上半环正电、下半环负电，对环心场均向下；净电荷为零不代表场为零。
\sourceproblem{6-T2}{23}{27}
\problemtext 无限均匀带电平面面密度 $\sigma$，距平面 $z>0$ 的P点，证明半径 $\sqrt3 z$ 的同轴圆盘区域贡献一半场强。
\workings 半径 $s$ 的环元电荷 $2\pi\sigma s\,\mathrm ds$，轴向场为
\[
E_z(a)=\frac{\sigma z}{2\varepsilon_0}\int_0^a\frac{s\,\mathrm ds}{(s^2+z^2)^{3/2}}
=\frac{\sigma}{2\varepsilon_0}\left(1-\frac{z}{\sqrt{z^2+a^2}}\right).
\]
当 $a=\sqrt3z$ 时括号等于 $1/2$，故圆盘贡献为无限平面场 $\sigma/(2\varepsilon_0)$ 的一半。
\crosscheck $a\to\infty$ 恢复无限平面结果；$a\to0$ 时场趋零。
\sourceproblem{6-T3}{24}{28}
\problemtext 半径 $R$ 的均匀带电半球曲面面密度 $\sigma$，求球心场强。
\workings 令 $\theta$ 为从半球向内对称轴（指向带电半球）量起的角，环元面积 $2\pi R^2\sin\theta\,\mathrm d\theta$。在球心，轴向场沿反向，
\[
E_{\parallel}=-\frac{\sigma}{2\varepsilon_0}\int_0^{\pi/2}\cos\theta\sin\theta\,\mathrm d\theta
=-\frac{\sigma}{4\varepsilon_0}.
\]
依原图，带电半球位于左方，若 $\sigma>0$，场指向右方开口，大小 $|\sigma|/(4\varepsilon_0)$。
\crosscheck 横向分量由轴对称抵消；两相反半球均匀同密度合成完整球面时中心场相消。
\sourceproblem{6-T4}{24}{28}
\problemtext 厚度 $d$ 的非导体无限平板，均匀体电荷密度 $\rho$，求板内外场强。
\workings 取中面为 $z=0$，垂直板面的正向为 $\bm e_z$，用对称薄盒高斯面得
\[
\bm E(z)=\begin{cases}
\rho z\,\bm e_z/\varepsilon_0,& |z|\le d/2,\\
\rho d\,\mathrm{sgn}(z)\bm e_z/(2\varepsilon_0),& |z|>d/2.
\end{cases}
\]
\crosscheck 两表面场连续，因为没有附加面电荷；中面场为零。原书未给横向尺寸，采用无限平板模型，有限板需考虑边缘效应。
\sourceproblem{6-T5}{25}{29}
\problemtext 点电荷 $q$ 位于正方体中心或一个顶点时，分别求每面电通量。
\workings 中心情形由六面等价及高斯定理，每面为 $q/(6\varepsilon_0)$。顶点情形，将八个同样正方体拼成大立方体使电荷居中心；每个大外面通量 $q/(6\varepsilon_0)$，分为四个等价小面，故原立方体不经过电荷的三个外面各为 $q/(24\varepsilon_0)$。经过该顶点的三个面除奇点外场平行于面，通量按去除奇点的积分为零。
\crosscheck 三个外面之和为 $q/(8\varepsilon_0)$；电荷位于边界，不能未经说明直接套用闭合面内部点电荷公式。
\sourceproblem{6-T6}{25}{29}
\problemtext 半径 $R$ 的球内 $\rho(r)=\rho_0(1-r/R)$，求内外场及最大场强位置。
\workings 球对称积分 $Q(r)=4\pi\rho_0(r^3/3-r^4/(4R))$（$r\le R$），总电荷 $Q(R)=\pi\rho_0R^3/3$，因此
\[
E_r=\begin{cases}
\frac{\rho_0}{\varepsilon_0}(r/3-r^2/(4R)),&0\le r\le R,\\
\frac{\rho_0 R^3}{12\varepsilon_0r^2},&r>R .
\end{cases}
\]
对内区大小求极值，$1/3-r/(2R)=0$ 给出 $r=2R/3$，$E_{\max}=|\rho_0|R/(9\varepsilon_0)$（$\rho_0\ne0$）。
\crosscheck 边界两式均给 $\rho_0R/(12\varepsilon_0)$，小于峰值；外区随 $r$ 单调减小，故为全局最大。
\sourceproblem{6-T7}{26}{30}
\problemtext 无限长均匀带电实心圆柱半径 $R$、体密度 $\rho$，证明内场并求外场。
\workings 取半径 $r$、长度 $L$ 的同轴高斯圆柱，端面通量为零。由 $2\pi rLE_r=\rho\pi L\min(r^2,R^2)/\varepsilon_0$ 得
\[
E_r=\begin{cases}
\rho r/(2\varepsilon_0),&0\le r\le R,\\
\rho R^2/(2\varepsilon_0r),&r>R .
\end{cases}
\]
\crosscheck 在 $R$ 处连续；外部相当于线密度 $\lambda=\pi R^2\rho$ 的无限线电荷。
\sourceproblem{6-T8}{26}{30}
\problemtext 同心球壳区域 $a<r<b$ 内体密度 $\rho=A/r$，空腔中心有点电荷 $Q$，求使带电区域场强大小恒定的 $A$。
\workings 对 $a<r<b$，包围净电荷为 $Q+4\pi A\int_a^r s\,\mathrm ds=Q+2\pi A(r^2-a^2)$。因此
\[
E_r=\frac{A}{2\varepsilon_0}+\frac{Q-2\pi Aa^2}{4\pi\varepsilon_0r^2}.
\]
要使连续区间内场强大小恒定，第二项系数必须为零，故 $A=Q/(2\pi a^2)$，恒定大小 $|Q|/(4\pi\varepsilon_0a^2)$。
\crosscheck 场强有向分量连续；非零恒定大小意味着该连通区间内方向符号不变，故消去 $r^{-2}$ 项是必要且充分的。量纲 $A$ 为电荷/面积。
